\section{The Crisis of 1976}

Roman authorities warned Lefebvre not to proceed with priestly ordinations planned for 29 June 1976. He ordained the candidates at Écône. The resulting measures unfolded in stages that are often compressed. A Holy See statement of 1 July said that the ordinations had incurred an automatic one-year suspension \emph{a collatione ordinum}, from conferring orders. A formal monition was dated 6 July and received on 11 July. On 22 July Lefebvre was notified of suspension \emph{a divinis}, barring the exercise of orders. SSPX sources preserve these records while disputing their justice and, in some arguments, their legal force. The public fact is less ambiguous: Lefebvre continued sacramental ministry despite the prohibitions.

The Mass at Lille on 29 August transformed a canonical confrontation into an international public event. It displayed both sides of the Society's appeal: the familiar older rite attracted Catholics disoriented by rapid reform, while Lefebvre's public address carried his critique of the Council and postconciliar reform beyond Écône. The ability to assemble supporters outside the seminary showed that suppression had not dissolved the movement.

Paul VI received Lefebvre on 11 September. In his Latin letter of 11 October the pope restated the central issues: rejection of the Council's authority, resistance to legitimate liturgical reform, disobedient ordinations, and the claim to act as guardian of tradition against the pope and bishops. The letter called Lefebvre to submission and offered a path of reconciliation, but it did not accept the continued seminary as the price of dialogue. Lefebvre sought practical guarantees for the older Mass and priestly work; the pope insisted that obedience and acceptance of the Council could not be bracketed.

\evidenceframe{Paul VI's public remarks of 1 September and letter of 11 October 1976, read with the dated suspension notifications preserved by the Society.}{The papal case joined council, liturgy, governance, and ministerial obedience; the penalties followed ordinations conducted after explicit prohibition.}{The pope's documents state the Holy See's judgment. The Society's necessity defense is a participant argument and is not silently converted into a recognized canonical exemption.}

\subsection{A pattern established}

The crisis fixed a pattern that would recur. First, a threatened institutional act---ordination, later episcopal consecration---created a deadline. Second, correspondence continued, but the parties assigned different minimum meanings to reconciliation. Third, Lefebvre acted to preserve the Society's future. Fourth, penalties or declarations followed. Fifth, the Society treated them as unjust consequences of fidelity rather than reasons to stop. Negotiation and unilateral action were thus not successive eras; they became concurrent methods.

The pattern also explains why ``the old Mass was forbidden'' is an insufficient causal summary. The liturgical question was real and deeply consequential, but Paul VI's acts addressed unauthorized ministry and rejection of ecclesiastical authority as well. Conversely, ``mere disobedience'' does not explain why the movement endured. For adherents the older rite, traditional seminary formation, and doctrinal opposition formed a credible account of emergency. Historical explanation must show the strength of that account without treating conviction as jurisdiction.

After 1976 SSPX priests celebrated, preached, taught, and founded works without a recognized canonical mission and, for acts requiring them, ordinary faculties. That practice normalized an exceptional ecclesiology in daily life: the visible structures of Catholic ministry were reproduced, but the usual bonds of incardination, diocesan authorization, and delegated faculty were contested or absent. The later appeal to supplied jurisdiction did not arise only as an abstract canonical theory; it grew within this sustained pastoral fact.
